\documentclass{article}
\input{../../lib.sty}

\title{\texttt{CompositionMeasure for MultiDP}}
\author{\MichaelShoemateAuthor}
\date{}

\begin{document}
\maketitle

\contrib
Proves soundness of the implementation of \rustdoc{combinators/sequential_composition/trait}{CompositionMeasure} for \texttt{MultiDP}.

\section{Hoare Triple}
\subsection*{Precondition}
\subsubsection*{Compiler-verified}
The output measure is derived from the common capabilities of all child measures.
Only capabilities with a common composition theorem are retained.

\subsubsection*{Caller-verified}
None.

\subsection*{Postcondition}
\begin{theorem}
\texttt{composability} returns \texttt{Ok(out)} if the composition of a vector of
privacy guarantees is bounded above by \texttt{compose(d\_mids)} under the
reported adaptivity and composability. Otherwise it returns an error.
\end{theorem}

\begin{proof}
The composition implementation evaluates each representation only when every
child supplies that representation. Native zCDP parameters are composed by
adding $\rho_i$ and source deltas. Pure RDP parameters are composed pointwise
by adding their curves; pure zCDP is embedded into RDP using the exact
$\alpha\rho_i$ relation. Gaussian-DP parameters, when enabled, are composed by
$\mu^2 = \sum_i \mu_i^2$.

Each accumulator uses directed interval arithmetic, so its returned value is
an upper bound on the corresponding exact composition. Approximate profiles
and tradeoff functions are not composed without a representation-specific
theorem. If no common supported capability exists, construction fails rather
than advertising an unsupported composition. Therefore every representation
retained in the resulting \texttt{PrivacyGuarantee} is a valid bound for the
same composed privacy relation.
\end{proof}

\bibliographystyle{alpha}
\bibliography{ref}

\end{document}
